% target: sec/08-society.tex
\src{A}{S}Four harms are documented. To find studies that measure a harm of oversampling directly, we searched PubMed on 10 November 2026 with one query for each harm that the literature names (calibration, leakage, privacy, and fairness between groups; online appendix, Section~S6) and screened the <<harm_records>> distinct records on title and abstract under a written rule: a study counts when it compares a model trained on oversampled data with one trained without resampling, or, for leakage, resampling before the split with resampling of the training data only, on a measure of the harm. The screen was done by AI sessions with a single screener per record (Section~\ref{sec:ai}). The rule admits <<harm_include>> studies, and <<harm_unsure>> records need the full text. Of the <<harm_calib>> studies that measure calibration, <<harm_calib_harm>> find it worse after oversampling, <<harm_calib_noharm>> find no harm and <<harm_calib_mixed>> finds mixed effects; all <<harm_leak>> that measure leakage find estimates inflated by resampling before the split; of the <<harm_fair>> that compare a disparity between groups, <<harm_fair_noharm>> find no harm and <<harm_fair_mixed>> finds mixed effects; and none of the <<harm_priv_records>> records of the privacy query compares an attack on oversampled data with one on the original data. Another <<harm_calib_nocomp>> records report the calibration of a model trained on oversampled data without any such comparison.
